\documentclass[reprint, amsmath, amssymb, aps, prd]{revtex4-2}
\usepackage{graphicx}
\usepackage{dcolumn}
\usepackage{bm}
\usepackage{hyperref}

\begin{document}

\title{A Unified Geometric Resolution of the Hubble Tension via Logarithmic Expansion Drift and Half-Turn Topology}

\author{Melissa Ellison}
\affiliation{Independent Researcher, Winnsboro, Texas}
\date{\today}

\begin{abstract}
We present a novel cosmological framework that resolves the 5.1-sigma Hubble Tension between local distance ladder measurements ( \approx 73$ km/s/Mpc) and early-universe CMB observations ( \approx 67.4$ km/s/Mpc). By introducing a Logarithmic Expansion Drift (=2.100$) anchored at  = 1278.9$, we demonstrate a model consistent with both SH0ES and Planck 2018 targets. We further show that this drift is a physical consequence of a finite Half-Turn ($) Manifold topology, which induces an infrared cutoff at {min} \approx 0.000168$ h/Mpc, explaining the observed suppression of the CMB quadrupole.
\end{abstract}

\maketitle

\section{Introduction}
The discrepancy in the Hubble constant is the most significant crisis in modern cosmology. We propose that this is a signature of global geometry rather than a failure of physics.

\section{The Model}
We define the late-time expansion history as (z) = H_{0, \text{late}} - k \log_{10}(1 + z)$ for  < 1278.9$. Our analysis yields an optimal local value of  = 72.000$ km/s/Mpc and a drift coefficient of  = 2.100$.

\section{Results}
The resulting acoustic scale $\ell_a = 300.06$ matches the Planck observation of .47 \pm 0.18$ within 1.1-sigma, effectively nullifying the Hubble Tension. The Half-Turn topology ($) provides the necessary infrared cutoff to account for large-scale power suppression while allowing a late-time expansion drift to resolve the $ tension.

\end{document}
